% Chapter 1: What is Control?
% IEEE-formatted LaTeX chapter for Control Systems Textbook
% Self-contained chapter - can be \input{} into main document

\chapter{What is Control?}
\label{ch:what_is_control}


%=============================================================================
\section{Historical Motivation: The Birth of Automatic Control}
\label{sec:history}
%=============================================================================

The science of automatic control did not emerge from abstract mathematical inquiry. Rather, it arose from urgent practical necessity---the need to regulate physical processes that humans could not monitor and adjust continuously. Understanding these historical origins provides essential intuition for why feedback mechanisms are structured as they are.

\subsection{The Ktesibios Water Clock (c. 270 BC)}
\label{subsec:ktesibios}

\begin{historicalbox}[The First Feedback Controller]
The earliest known feedback control device was invented by Ctesibius of Alexandria (also written Ktesibios), a Greek inventor and mathematician working in Ptolemaic Egypt around 270 BC. His ingenious water clock, or \textit{clepsydra}, solved a problem that had plagued timekeeping for millennia. Ktesibios was the son of a barber and initially followed his father's trade before becoming one of the most prolific inventors of the ancient world. His contributions extended beyond the water clock to include the pipe organ (hydraulis), force pump, and various pneumatic devices. The float valve mechanism he invented for the clepsydra represents humanity's first documented automatic control system---a remarkable achievement that predates the Industrial Revolution by over two millennia.
\end{historicalbox}

The earliest known feedback control device was invented by Ctesibius of Alexandria (also written Ktesibios), a Greek inventor and mathematician working in Ptolemaic Egypt around 270 BC. His ingenious water clock, or \textit{clepsydra}, solved a problem that had plagued timekeeping for millennia.

\subsubsection{The Problem: Variable Flow Rate}

Prior water clocks operated on a simple principle: water draining from a vessel through an orifice at the bottom would indicate time by the falling water level. However, these devices suffered from a fundamental flaw rooted in fluid dynamics. The flow rate $Q$ through an orifice depends on the pressure head $h$ above it, following Torricelli's law:
\begin{equation}
Q = C_d A \sqrt{2gh}
\label{eq:torricelli}
\end{equation}
where $C_d$ is the discharge coefficient, $A$ is the orifice area, and $g$ is gravitational acceleration. As water drained and $h$ decreased, the flow rate diminished, causing the clock to run progressively slower. A clock accurate in the morning would lose significant time by evening.

\subsubsection{The Solution: Float Valve Regulation}

Ktesibios introduced a remarkable solution: a separate reservoir maintained at constant level by a float valve, which then fed the timing vessel. Water entered the regulating reservoir from an external source at a variable rate. A float, constructed from cork or a hollow sphere, rose and fell with the water level. This float was mechanically connected to a conical valve known as the \textit{cock}. As the water level rose, the float lifted and progressively closed the inlet valve, restricting flow. Conversely, as the water level fell, the float dropped and opened the valve to admit more water.

This mechanism maintained a nearly constant head $h$ in the regulating reservoir, ensuring constant flow rate to the timing vessel regardless of variations in the external water supply. The float valve constitutes a true feedback controller: the controlled variable (water level) directly influences the manipulated variable (inlet valve position) through a mechanical linkage.

\subsubsection{Why Open-Loop Failed}

Consider the alternative: manually setting a fixed inlet flow rate intended to match the outlet flow. Any disturbance---variation in source pressure, partial clogging of the inlet, evaporation, temperature changes affecting viscosity---would cause the level to drift. Without continuous human intervention, accuracy was impossible. The float valve made the system \textit{self-regulating}, automatically compensating for disturbances that the designer could not anticipate or eliminate.

\subsection{Cornelis Drebbel's Thermostat (c. 1620)}
\label{subsec:drebbel}

The Dutch inventor Cornelis Drebbel constructed what is considered the first temperature feedback control system around 1620. His application was the regulation of furnace temperature for alchemical operations and, notably, for chicken egg incubators---where maintaining precise temperature was essential for successful hatching.

\subsubsection{The Problem: Temperature Fluctuation}

Maintaining constant temperature in a furnace or incubator using only fuel rate adjustment is extraordinarily difficult. Heat loss varies with ambient temperature, fuel quality varies, and the thermal mass of the system creates delays between fuel input and temperature response. An operator attempting to maintain temperature by periodic adjustment inevitably produces oscillations: overcorrection leads to overheating, followed by under-fueling and cooling, in an endless cycle.

\subsubsection{The Mechanism}

Drebbel's thermostat employed a mercury-and-alcohol thermometer connected to a damper controlling air flow to the furnace. A glass vessel containing alcohol was placed inside the incubator, where the alcohol expanded or contracted in response to temperature changes. This expansion or contraction moved a column of mercury in a connected tube, and the mercury column was mechanically linked to a damper or flap controlling airflow. When temperature increased, the alcohol expanded, causing the mercury to rise and the damper to close, thereby reducing combustion. When temperature decreased, the alcohol contracted, the mercury fell, and the damper opened to increase combustion.

This represents a proportional feedback controller, where the damper position varied continuously with temperature deviation from the setpoint (determined by the mechanical linkage geometry).

\subsection{James Watt's Centrifugal Governor (1788)}
\label{subsec:watt}

The device that launched control theory as a mathematical discipline was James Watt's centrifugal governor for steam engines, patented in 1788. While Watt did not invent the centrifugal mechanism (similar devices regulated windmill speeds), his application to steam engines created the technological foundation of the Industrial Revolution---and the theoretical problems that birthed modern control theory.

\subsubsection{The Problem: Catastrophic Speed Variation}

Early steam engines suffered from dangerous speed fluctuations. The fundamental issue arose from the mismatch between steam power input and mechanical load. Industrial loads such as pumps, looms, and mills varied moment to moment; a loom engaging produced sudden load while thread breaking released it. Steam pressure variations compounded the problem, as boiler output fluctuated with fuel quality, water level, and draft conditions. Furthermore, some loads exhibited positive feedback instabilities, with decreasing resistance at higher speeds creating potential runaway conditions.

Without automatic regulation, a sudden load decrease could cause the engine to ``run away,'' potentially to mechanical destruction. Conversely, sudden load increases could stall the engine. Human operators could not react quickly enough, and continuous manual throttle adjustment was impractical in factory settings.

\subsubsection{The Centrifugal Governor Mechanism}

Watt's governor operated on an elegant principle relating rotational speed to centrifugal force. Two heavy balls were mounted on hinged arms connected to a vertical spindle, which was driven by the engine through gearing and rotated proportionally to engine speed. As speed increased, centrifugal force swung the balls outward and upward, and the ball arm geometry converted this outward motion to vertical motion of a sliding collar. The collar was linked to the steam throttle valve, creating the feedback mechanism: higher speed caused the balls to rise, raising the collar and closing the throttle, which admitted less steam and decreased speed. Conversely, lower speed caused the balls to fall, lowering the collar and opening the throttle, which admitted more steam and increased speed.

The equilibrium condition occurs when centrifugal force exactly balances gravity and spring forces, determining a specific speed at which the throttle position stabilizes.

\subsubsection{Mathematical Analysis: Maxwell's 1868 Paper}

The Watt governor worked remarkably well in practice, but some installations exhibited puzzling oscillations---the engine speed would hunt continuously about the setpoint rather than settling. This phenomenon prompted James Clerk Maxwell to publish ``On Governors'' in 1868, founding the mathematical theory of feedback control.

Maxwell modeled the governor-engine system as a set of differential equations and analyzed their stability using linearization. He showed that oscillations arose when certain relationships between governor parameters (ball mass, arm length, spindle friction) and engine parameters (inertia, throttle sensitivity) were violated. The mathematical condition for stability---that all roots of the characteristic equation have negative real parts---was a precursor to the Routh-Hurwitz criterion developed later.

\subsubsection{The Fundamental Lesson}

The Watt governor illustrates the central principle motivating feedback control:

\begin{tcolorbox}[colback=blue!5!white,colframe=blue!75!black,title=The Feedback Principle,sharp corners]
Feedback control enables a system to regulate itself against disturbances that cannot be predicted, measured, or eliminated at their source. The controller need not ``know'' why the output deviated---only that it deviated---to take corrective action.
\end{tcolorbox}

An open-loop alternative---setting the throttle to a fixed position calculated to produce the desired speed---would fail immediately upon any load change, boiler pressure variation, or friction change. The governor succeeds because it continuously measures the actual output (speed) and adjusts the input (steam) accordingly.

%=============================================================================
\section{Modern Applications: Open-Loop vs.\ Closed-Loop Systems}
\label{sec:modern_applications}
%=============================================================================

The distinction between open-loop and closed-loop control pervades modern technology, though users rarely consider which paradigm governs the devices they use.

\subsection{Open-Loop Control Systems}

Open-loop systems apply a predetermined input without measuring the output. They rely on accurate system models and the absence of disturbances.

\subsubsection{Household Examples}
Common household appliances illustrate open-loop operation. A timer-based toaster applies heat for a fixed duration regardless of actual bread browning, and variations in bread moisture, thickness, and initial temperature cause inconsistent results. Similarly, many basic washing machines run fixed cycles assuming standard loads, so overloading or underloading produces suboptimal cleaning. Microwave ovens apply power for a user-specified time without measuring food temperature, though some modern units add temperature probes, thereby converting to closed-loop operation.

\subsubsection{Industrial Examples}
Industrial settings also employ open-loop control in certain applications. In low-precision stepper motor positioning, motors are commanded a specific number of steps assuming each step produces exact displacement, but missed steps accumulate as position error. Some industrial processes also add chemical reagents on a time schedule rather than measuring concentration, relying on predictable reaction kinetics.

\subsubsection{When Open-Loop Suffices}
Open-loop control is appropriate under specific conditions. The system model must be accurate and stable, and disturbances must be negligible or repeatable. Additionally, the cost of feedback sensing should exceed the value of improved performance, and the consequences of error must be acceptable. When these conditions are met, the simplicity and lower cost of open-loop control make it the preferred choice.

\subsection{Closed-Loop Control Systems}

Closed-loop systems measure the output and adjust the input to minimize error between the output and desired reference.

\subsubsection{Automotive Examples}
Modern automobiles employ numerous closed-loop control systems. Cruise control uses a speed sensor to continuously measure vehicle velocity, adjusting the throttle to maintain the setpoint despite hills, wind, and rolling resistance variations. Anti-lock braking systems (ABS) use wheel speed sensors to detect impending lockup and modulate brake pressure to maintain optimal slip ratio. Electronic stability control employs yaw rate and lateral acceleration sensors to detect skids, applying individual wheel braking to restore the intended trajectory.

\subsubsection{Industrial Process Control}
Industrial processes rely heavily on closed-loop control for safety and quality. In refinery operations, temperature, pressure, flow rate, and composition sensors feed controllers that adjust heaters, valves, and pumps. Power grid frequency control uses generator governors---descendants of Watt's device---along with automatic generation control to maintain 50/60 Hz despite load variations. Semiconductor fabrication requires sub-nanometer positioning, achieved through continuous feedback from interferometric sensors.

\subsubsection{Robotics and Automation}
Robotics and automation depend critically on feedback control. Robot arm positioning uses encoder feedback to achieve sub-millimeter repeatability despite payload variations. Drone flight control integrates data from accelerometers, gyroscopes, barometers, and GPS to maintain attitude and position. Self-driving vehicles employ LIDAR, cameras, and radar for perception, with control algorithms maintaining lane position and vehicle following distance.

\subsubsection{Biological Feedback Systems}
Biological organisms are extraordinarily sophisticated feedback control systems. Thermoregulation exemplifies this: the hypothalamus maintains core body temperature at 37°C through sweating, shivering, and vasodilation or vasoconstriction. Blood glucose regulation employs insulin and glucagon as a dual-controller system maintaining glucose homeostasis. The pupillary light reflex adjusts pupil diameter to maintain appropriate retinal illumination, while the vestibulo-ocular reflex compensates eye position for head rotation, thereby stabilizing the visual field.

%=============================================================================
\section{Mathematical Foundations}
\label{sec:math_foundations}
%=============================================================================

We now establish the mathematical framework for analyzing control systems, beginning with precise definitions and progressing through the fundamental equations of open-loop and closed-loop operation.

\subsection{System Definitions and Notation}

\begin{definition}[System]
A \textbf{system} is a collection of components that act together to perform a function. In control engineering, we model systems as operators that transform input signals into output signals.
\end{definition}

\begin{definition}[Signals]
We define the following signals in a control system. The \textbf{reference input} $r(t)$ represents the setpoint or desired value. The \textbf{control input} $u(t)$ is the manipulated variable or actuator command. The \textbf{output} $y(t)$ denotes the controlled variable or measured response. External influences enter the system as the \textbf{disturbance} $d(t)$, representing uncontrolled factors, and \textbf{noise} $n(t)$, representing measurement corruption. Finally, the \textbf{error signal} $e(t)$ is defined as the difference between reference and output: $e(t) = r(t) - y(t)$.
\end{definition}

\begin{definition}[Transfer Function]
For a linear time-invariant (LTI) system, the \textbf{transfer function} $G(s)$ relates the Laplace transform of the output to the Laplace transform of the input:
\begin{equation}
G(s) = \frac{Y(s)}{U(s)}
\label{eq:transfer_function}
\end{equation}
where $Y(s) = \mathcal{L}\{y(t)\}$ and $U(s) = \mathcal{L}\{u(t)\}$.
\end{definition}

\subsection{Block Diagram Representation}

Control systems are represented using block diagrams, a graphical notation with several standard elements. Blocks represent transfer functions, with input entering from the left and output exiting to the right. Arrows indicate signal flow direction. Summing junctions, depicted as circles with $\Sigma$ or $+/-$ symbols, add or subtract signals. Pickoff points, shown as filled circles, duplicate a signal for routing to multiple destinations.

% Open-loop block diagram
\begin{figure}[htbp]
\centering
\begin{tikzpicture}[auto, node distance=2.5cm, >=latex',
    block/.style={rectangle, draw, minimum height=2em, minimum width=3em},
    sum/.style={circle, draw, inner sep=0pt, minimum size=6mm},
    input/.style={coordinate},
    output/.style={coordinate}]
    % Blocks
    \node [input, name=rinput] (rinput) {};
    \node [block, right of=rinput] (controller) {$K$};
    \node [sum, right of=controller, node distance=2cm] (sum1) {};
    \node [block, right of=sum1, node distance=2cm] (plant) {$G(s)$};
    \node [output, right of=plant] (output) {};

    % Disturbance input
    \node [input, above of=sum1, node distance=1.5cm] (disturbance) {};

    % Connections
    \draw [->] (rinput) -- node {$R(s)$} (controller);
    \draw [->] (controller) -- node {$U(s)$} (sum1);
    \draw [->] (sum1) -- (plant);
    \draw [->] (plant) -- node {$Y(s)$} (output);
    \draw [->] (disturbance) -- node {$D(s)$} (sum1);

    % Sum labels
    \node at (sum1.south east) {\small $+$};
    \node at (sum1.north) {\small $+$};
\end{tikzpicture}
\caption{Open-loop control system block diagram. The reference input $R(s)$ passes through the controller $K$ to produce the control signal $U(s)$, which combines with the disturbance $D(s)$ at the plant input. No feedback path exists.}
\label{fig:openloop_block}
\end{figure}

\subsection{Open-Loop Control Analysis}

In open-loop control, the controller operates on the reference input without knowledge of the output. The control input is:
\begin{equation}
U(s) = K \cdot R(s)
\label{eq:openloop_control}
\end{equation}
where $K$ represents the controller (possibly a constant gain, possibly a transfer function).

The output, including disturbance effects, is:
\begin{equation}
Y(s) = G(s) \cdot U(s) + G(s) \cdot D(s) = G(s) K \cdot R(s) + G(s) \cdot D(s)
\label{eq:openloop_output}
\end{equation}

For ideal tracking with no disturbance ($D(s) = 0$), we desire $Y(s) = R(s)$, requiring:
\begin{equation}
G(s) K = 1 \quad \Rightarrow \quad K = \frac{1}{G(s)} = G^{-1}(s)
\label{eq:openloop_ideal}
\end{equation}

This inverse controller approach has severe limitations. The plant $G(s)$ must be exactly known, and the plant must be minimum-phase (having no right-half-plane zeros) for stable inversion. Any model error $\Delta G$ causes proportional output error, and disturbances $D(s)$ appear directly at the output, amplified by $G(s)$.

\subsection{Closed-Loop Control Analysis}

In closed-loop control, the error between reference and measured output drives the controller:

% Closed-loop block diagram
\begin{figure}[htbp]
\centering
\begin{tikzpicture}[auto, node distance=2cm, >=latex',
    block/.style={rectangle, draw, minimum height=2em, minimum width=3em},
    sum/.style={circle, draw, inner sep=0pt, minimum size=6mm},
    input/.style={coordinate},
    output/.style={coordinate}]
    % Blocks
    \node [input] (rinput) {};
    \node [sum, right of=rinput] (sum1) {};
    \node [block, right of=sum1] (controller) {$K(s)$};
    \node [sum, right of=controller] (sum2) {};
    \node [block, right of=sum2] (plant) {$G(s)$};
    \node [output, right of=plant, node distance=2.5cm] (output) {};

    % Feedback path
    \node [block, below of=plant, node distance=1.5cm] (sensor) {$H(s)$};

    % Disturbance
    \node [input, above of=sum2, node distance=1.5cm] (dist) {};

    % Forward path connections
    \draw [->] (rinput) -- node[pos=0.2] {$R(s)$} (sum1);
    \draw [->] (sum1) -- node {$E(s)$} (controller);
    \draw [->] (controller) -- node {$U(s)$} (sum2);
    \draw [->] (sum2) -- (plant);
    \draw [->] (plant) -- node [name=y, pos=0.7] {$Y(s)$} (output);
    \draw [->] (dist) -- node {$D(s)$} (sum2);

    % Feedback path
    \draw [->] (y) |- (sensor);
    \draw [->] (sensor) -| node[pos=0.9] {\small $-$} (sum1);

    % Sum labels
    \node at (sum1.north east) {\small $+$};
    \node at (sum2.south east) {\small $+$};
    \node at (sum2.north) {\small $+$};
\end{tikzpicture}
\caption{Closed-loop control system block diagram. The error $E(s) = R(s) - H(s)Y(s)$ drives the controller, and the feedback path through sensor $H(s)$ enables automatic disturbance rejection.}
\label{fig:closedloop_block}
\end{figure}

The mathematical relationships are:
\begin{align}
E(s) &= R(s) - H(s) Y(s) \label{eq:error_def}\\
U(s) &= K(s) E(s) \label{eq:controller_action}\\
Y(s) &= G(s) \left[ U(s) + D(s) \right] \label{eq:plant_response}
\end{align}

Substituting (\ref{eq:error_def}) into (\ref{eq:controller_action}) and then into (\ref{eq:plant_response}):
\begin{equation}
Y(s) = G(s) \left[ K(s) \left( R(s) - H(s) Y(s) \right) + D(s) \right]
\end{equation}

Expanding and collecting terms:
\begin{equation}
Y(s) + G(s) K(s) H(s) Y(s) = G(s) K(s) R(s) + G(s) D(s)
\end{equation}
\begin{equation}
Y(s) \left[ 1 + G(s) K(s) H(s) \right] = G(s) K(s) R(s) + G(s) D(s)
\end{equation}

Solving for $Y(s)$:
\begin{equation}
\boxed{Y(s) = \frac{G(s) K(s)}{1 + G(s) K(s) H(s)} R(s) + \frac{G(s)}{1 + G(s) K(s) H(s)} D(s)}
\label{eq:closedloop_general}
\end{equation}

Define the \textbf{loop transfer function}:
\begin{equation}
L(s) \triangleq G(s) K(s) H(s)
\label{eq:loop_tf}
\end{equation}

Then the closed-loop transfer functions are:

\begin{tcolorbox}[colback=green!5!white,colframe=green!75!black,title=Fundamental Closed-Loop Transfer Functions,sharp corners]
\textbf{Reference to Output (Complementary Sensitivity):}
\begin{equation}
T(s) = \frac{Y(s)}{R(s)} = \frac{G(s) K(s)}{1 + L(s)}
\label{eq:complementary_sensitivity}
\end{equation}

\textbf{Disturbance to Output:}
\begin{equation}
\frac{Y(s)}{D(s)} = \frac{G(s)}{1 + L(s)}
\label{eq:disturbance_tf}
\end{equation}

\textbf{Reference to Error (Sensitivity):}
\begin{equation}
S(s) = \frac{E(s)}{R(s)} = \frac{1}{1 + L(s)}
\label{eq:sensitivity}
\end{equation}
\end{tcolorbox}

For unity feedback ($H(s) = 1$), these simplify further. Note the fundamental relationship:
\begin{equation}
S(s) + T(s) = 1
\label{eq:sensitivity_complementary}
\end{equation}

\subsection{Why Feedback Reduces Sensitivity}

Consider a plant with nominal transfer function $G_0(s)$ subject to multiplicative uncertainty:
\begin{equation}
G(s) = G_0(s)(1 + \Delta(s))
\label{eq:uncertain_plant}
\end{equation}
where $\Delta(s)$ represents model error or parameter variation.

\subsubsection{Open-Loop Sensitivity}

For open-loop control with $K = 1/G_0$ designed for the nominal plant:
\begin{equation}
Y = G \cdot K \cdot R = G_0(1+\Delta) \cdot \frac{1}{G_0} \cdot R = (1+\Delta) R
\end{equation}

The output error due to plant variation is:
\begin{equation}
\Delta Y_{\text{OL}} = \Delta \cdot R
\label{eq:openloop_sensitivity}
\end{equation}

The open-loop sensitivity to plant variations equals 1---any plant change appears directly at the output.

\subsubsection{Closed-Loop Sensitivity}

For closed-loop control with high loop gain $L_0 = G_0 K H \gg 1$:
\begin{equation}
T = \frac{G K}{1 + G K H} = \frac{G_0(1+\Delta)K}{1 + G_0(1+\Delta)K H}
\end{equation}

For large $L_0$:
\begin{equation}
T \approx \frac{G_0(1+\Delta)K}{G_0(1+\Delta)K H} = \frac{1}{H}
\end{equation}

The closed-loop transfer function becomes independent of the plant $G$, depending only on the sensor $H$! More precisely, the sensitivity of $T$ to plant variations is:
\begin{equation}
\frac{\partial T}{\partial G} \cdot \frac{G}{T} = \frac{1}{1 + L} = S
\label{eq:closedloop_sensitivity}
\end{equation}

The sensitivity function $S(s)$ quantifies how much plant variations affect the closed-loop response. At frequencies where $|L(j\omega)| \gg 1$, we have $|S(j\omega)| \ll 1$, and the system is insensitive to plant variations.

\begin{theorem}[Sensitivity Reduction]
Feedback reduces sensitivity to plant variations by the factor $(1 + L)^{-1}$. At frequencies where the loop gain magnitude $|L(j\omega)| = 100$ (40 dB), plant variations are attenuated by a factor of 100 at the output.
\end{theorem}

\subsection{Steady-State Error Analysis}

For asymptotic tracking of reference inputs, we analyze the steady-state error using the Final Value Theorem. For a stable closed-loop system:
\begin{equation}
e_{ss} = \lim_{t \to \infty} e(t) = \lim_{s \to 0} s \cdot E(s) = \lim_{s \to 0} s \cdot S(s) \cdot R(s)
\label{eq:steady_state_error}
\end{equation}

\subsubsection{System Type and Error Constants}

The \textbf{system type} $n$ is defined as the number of free integrators in the loop transfer function:
\begin{equation}
L(s) = \frac{K_L \prod_i (s + z_i)}{s^n \prod_j (s + p_j)}
\label{eq:system_type}
\end{equation}

The error constants are defined as:
\begin{align}
K_p &= \lim_{s \to 0} L(s) && \text{(position constant)} \label{eq:kp}\\
K_v &= \lim_{s \to 0} s \cdot L(s) && \text{(velocity constant)} \label{eq:kv}\\
K_a &= \lim_{s \to 0} s^2 \cdot L(s) && \text{(acceleration constant)} \label{eq:ka}
\end{align}

\begin{table}[htbp]
\centering
\caption{Steady-State Error for Different System Types}
\label{tab:steady_state_error}
\begin{tabular}{cccc}
\toprule
\textbf{System Type} & \textbf{Step Input} $R(s) = 1/s$ & \textbf{Ramp Input} $R(s) = 1/s^2$ & \textbf{Parabolic} $R(s) = 1/s^3$ \\
\midrule
Type 0 & $\dfrac{1}{1+K_p}$ & $\infty$ & $\infty$ \\[2ex]
\midrule
Type 1 & 0 & $\dfrac{1}{K_v}$ & $\infty$ \\[2ex]
\midrule
Type 2 & 0 & 0 & $\dfrac{1}{K_a}$ \\[2ex]
\bottomrule
\end{tabular}
\end{table}

%=============================================================================
\section{Practical Laboratory: DC Motor Position Control}
\label{sec:lab}
%=============================================================================

This laboratory exercise demonstrates the fundamental difference between open-loop and closed-loop control using readily available components. Students will implement both control strategies on a DC motor position system and quantify the performance improvement achieved through feedback.

\subsection{Learning Objectives}

Upon completion of this laboratory, students will be able to explain why open-loop control fails under disturbances and model uncertainty, implement a basic closed-loop position controller using encoder feedback, measure and compare steady-state error, disturbance rejection, and repeatability, and tune a proportional controller gain while observing its effects on system behavior.

\subsection{Required Components}

\begin{table}[htbp]
\centering
\caption{Bill of Materials for Motor Control Laboratory}
\label{tab:bom}
\begin{tabular}{lll}
\toprule
\textbf{Component} & \textbf{Specification} & \textbf{Qty} \\
\midrule
Arduino Uno or Nano & ATmega328P microcontroller & 1 \\
DC Motor with Encoder & 12V, 200+ PPR quadrature encoder & 1 \\
Motor Driver & L298N or TB6612FNG H-bridge & 1 \\
Potentiometer & 10k$\Omega$ linear taper & 1 \\
Power Supply & 12V, 2A DC adapter & 1 \\
Breadboard & Full-size & 1 \\
Jumper Wires & Male-male, assorted & 20 \\
Position Indicator & Pointer or dial attached to motor shaft & 1 \\
Reference Scale & Printed protractor or degree markings & 1 \\
\bottomrule
\end{tabular}
\end{table}

For advanced experiments, optional components include a 3D-printed dial with degree markings for the motor shaft, a second potentiometer for setpoint adjustment, and serial plotting software such as Arduino Serial Plotter or Processing.

\subsection{Wiring Configuration}

\begin{table}[htbp]
\centering
\caption{Motor Control System Wiring Connections}
\label{tab:wiring}
\begin{tabular}{lll}
\toprule
\textbf{Category} & \textbf{Source} & \textbf{Destination} \\
\midrule
\multirow{4}{*}{Power} & 12V Supply & Motor Driver VCC \\
& Motor Driver GND & Common Ground / Arduino GND \\
& Arduino 5V & Potentiometer Terminal 1 \\
& Arduino GND & Potentiometer Terminal 2 \\
\midrule
\multirow{3}{*}{Motor Driver} & Driver IN1 & Arduino D5 (PWM) \\
& Driver IN2 & Arduino D6 (PWM) \\
& Driver ENA & Arduino D9 (PWM) or 5V jumper \\
\midrule
\multirow{4}{*}{Encoder} & Encoder VCC & Arduino 5V \\
& Encoder GND & Arduino GND \\
& Encoder A & Arduino D2 (interrupt capable) \\
& Encoder B & Arduino D3 (interrupt capable) \\
\midrule
Potentiometer & Wiper & Arduino A0 \\
\bottomrule
\end{tabular}
\end{table}

\subsection{Software Implementation}

\subsubsection{Part 1: Encoder Reading (Common to Both Experiments)}

\begin{algorithm}
\caption{Quadrature Encoder Reading with Interrupts}
\label{alg:encoder}
\begin{algorithmic}[1]
\State \textbf{Initialize:}
\State \hspace{1em} $\textit{encoderCount} \gets 0$
\State \hspace{1em} $\textit{lastEncoded} \gets 0$
\State \hspace{1em} Configure encoder pins A and B as inputs with pull-up resistors
\State \hspace{1em} Attach interrupt handlers to both encoder channels
\State
\Procedure{UpdateEncoder}{} \Comment{Called on any encoder pin change}
    \State $\textit{MSB} \gets$ read encoder channel A
    \State $\textit{LSB} \gets$ read encoder channel B
    \State $\textit{encoded} \gets (\textit{MSB} \times 2) + \textit{LSB}$
    \State $\textit{sum} \gets (\textit{lastEncoded} \times 4) + \textit{encoded}$
    \If{$\textit{sum}$ indicates forward rotation}
        \State $\textit{encoderCount} \gets \textit{encoderCount} + 1$
    \ElsIf{$\textit{sum}$ indicates reverse rotation}
        \State $\textit{encoderCount} \gets \textit{encoderCount} - 1$
    \EndIf
    \State $\textit{lastEncoded} \gets \textit{encoded}$
\EndProcedure
\State
\Function{GetPositionDegrees}{}
    \State $\textit{COUNTS\_PER\_REV} \gets 400$ \Comment{Adjust for encoder PPR}
    \State \Return $(\textit{encoderCount} / \textit{COUNTS\_PER\_REV}) \times 360$
\EndFunction
\end{algorithmic}
\end{algorithm}

\subsubsection{Part 2: Open-Loop Control Experiment}

\begin{algorithm}
\caption{Open-Loop Position Control (Timed Pulses)}
\label{alg:openloop}
\begin{algorithmic}[1]
\State \textbf{Constants} (students must calibrate):
\State \hspace{1em} $\textit{CAL\_PWM} \gets 150$ \Comment{Motor speed setting}
\State \hspace{1em} $\textit{CAL\_TIME\_90DEG} \gets 500$ \Comment{Milliseconds for 90 degrees}
\State
\State \textbf{Initialize:}
\State \hspace{1em} Configure motor pins as outputs
\State \hspace{1em} Initialize encoder (see Algorithm~\ref{alg:encoder})
\State
\Procedure{SetMotorPWM}{$\textit{pwm}$}
    \If{$\textit{pwm} > 0$}
        \State Set forward direction with magnitude $\textit{pwm}$
    \ElsIf{$\textit{pwm} < 0$}
        \State Set reverse direction with magnitude $|\textit{pwm}|$
    \Else
        \State Stop motor
    \EndIf
\EndProcedure
\State
\Procedure{MoveOpenLoop}{$\textit{targetDegrees}$}
    \State $\textit{timeRequired} \gets |\textit{targetDegrees}| \times (\textit{CAL\_TIME\_90DEG} / 90)$
    \State $\textit{direction} \gets \text{sign}(\textit{targetDegrees})$
    \State $\textit{startPos} \gets \text{GetPositionDegrees}()$
    \State \Call{SetMotorPWM}{$\textit{direction} \times \textit{CAL\_PWM}$}
    \State Wait for $\textit{timeRequired}$ milliseconds
    \State \Call{SetMotorPWM}{$0$}
    \State Wait 500 ms for motor to settle
    \State $\textit{endPos} \gets \text{GetPositionDegrees}()$
    \State $\textit{error} \gets \textit{targetDegrees} - (\textit{endPos} - \textit{startPos})$
    \State Report target, actual position, and error
\EndProcedure
\State
\State \textbf{Main Loop:}
\While{true}
    \If{user input available}
        \State $\textit{target} \gets$ read target angle from user
        \State Reset encoder count to zero
        \State \Call{MoveOpenLoop}{$\textit{target}$}
    \EndIf
\EndWhile
\end{algorithmic}
\end{algorithm}

\subsubsection{Part 3: Closed-Loop Control Experiment}

\begin{algorithm}
\caption{Closed-Loop Proportional Position Control}
\label{alg:closedloop}
\begin{algorithmic}[1]
\State \textbf{Parameters} (students tune these):
\State \hspace{1em} $K_p \gets 2.0$ \Comment{Proportional gain}
\State \hspace{1em} $\textit{DEADBAND} \gets 1.0$ \Comment{Degrees; stop when error below this}
\State \hspace{1em} $\textit{MAX\_PWM} \gets 200$
\State \hspace{1em} $\textit{MIN\_PWM} \gets 30$ \Comment{Minimum to overcome friction}
\State
\State \textbf{State Variables:}
\State \hspace{1em} $\textit{setpoint} \gets 0$
\State
\State \textbf{Initialize:}
\State \hspace{1em} Configure encoder and motor pins
\State
\State \textbf{Main Control Loop} (runs at 100 Hz):
\While{true}
    \If{user command received}
        \If{command is `S'}
            \State $\textit{setpoint} \gets$ new setpoint value
        \ElsIf{command is `K'}
            \State $K_p \gets$ new gain value
        \EndIf
    \EndIf
    \State
    \State $\textit{position} \gets \text{GetPositionDegrees}()$
    \State $\textit{error} \gets \textit{setpoint} - \textit{position}$
    \State $\textit{pwmOutput} \gets 0$
    \State
    \If{$|\textit{error}| > \textit{DEADBAND}$}
        \State $\textit{pwmOutput} \gets K_p \times \textit{error}$ \Comment{Proportional control}
        \State $\textit{pwmOutput} \gets \text{clamp}(\textit{pwmOutput}, -\textit{MAX\_PWM}, \textit{MAX\_PWM})$
        \If{$|\textit{pwmOutput}| < \textit{MIN\_PWM}$}
            \State $\textit{pwmOutput} \gets \text{sign}(\textit{error}) \times \textit{MIN\_PWM}$
        \EndIf
    \EndIf
    \State
    \State \Call{SetMotorPWM}{$\textit{pwmOutput}$}
    \State Report setpoint, position, error, and PWM (every 100 ms)
    \State Wait 10 ms
\EndWhile
\end{algorithmic}
\end{algorithm}

\subsection{Experimental Procedure}

\subsubsection{Experiment A: Open-Loop Calibration and Testing}

Begin by uploading the open-loop code to the Arduino. For calibration, command 90-degree movements and adjust the \texttt{CAL\_TIME\_90DEG} parameter until the motor consistently moves 90 degrees, then record this calibrated value. Next, conduct a repeatability test by commanding ten 90-degree movements, recording the actual positions achieved, and calculating the mean and standard deviation of the error. For the disturbance test, while the motor is at rest, manually rotate the shaft, then command a 90-degree movement and observe that the system has no knowledge of this disturbance. Finally, perform a load test by attaching a small mass to the shaft and repeating the 90-degree movements; observe the increased error due to the changed dynamics.

\subsubsection{Experiment B: Closed-Loop Control}

Upload the closed-loop code to the Arduino and start with $K_p = 1.0$, commanding position changes and observing the response. For gain tuning, increase $K_p$ gradually and observe the system behavior: at low $K_p$ values, the response will be sluggish with large steady-state error; at medium $K_p$ values, the response becomes crisp with small error; at high $K_p$ values, oscillation or overshoot will appear. Test disturbance rejection by holding the motor at a position, manually pushing the shaft, and observing the automatic return to setpoint. Conduct a repeatability test by commanding ten 90-degree movements and comparing the error statistics to those from the open-loop experiment. Finally, perform a load test by attaching the same mass used in Experiment A and observing that the closed-loop system compensates automatically.

\subsection{Data Collection and Analysis}

Students should complete the following table:

\begin{table}[htbp]
\centering
\caption{Experimental Results Comparison}
\label{tab:results}
\begin{tabular}{lcc}
\toprule
\textbf{Metric} & \textbf{Open-Loop} & \textbf{Closed-Loop} \\
\midrule
Mean error (90° target) & & \\
Std. dev. of error & & \\
Maximum error observed & & \\
Time to reach 2° of target & & \\
Disturbance recovery & None & \\
Effect of added load & & \\
\bottomrule
\end{tabular}
\end{table}

\subsubsection{Discussion Questions}
Students should consider and discuss the following questions. First, why does open-loop control fail when load is added, while closed-loop adapts automatically? Second, what limits how high $K_p$ can be set, and what physical phenomena cause oscillation at high gain? Third, how would you modify the controller to eliminate steady-state error entirely? Fourth, what is the role of the deadband, and what happens if it is set too small or too large?

%=============================================================================
\section{Worked Examples}
\label{sec:examples}
%=============================================================================

\begin{examplebox}[Steady-State Error Comparison]
A DC motor position system has plant transfer function $G(s) = \dfrac{10}{s(s+2)}$. Compare the steady-state error to a unit ramp input for (a) open-loop control with $K = 1$, and (b) closed-loop control with proportional controller $K = 5$ and unity feedback.

\textbf{Solution:}
\end{examplebox}

\begin{example}[Steady-State Error Comparison]
\label{ex:sse}
A DC motor position system has plant transfer function $G(s) = \dfrac{10}{s(s+2)}$. Compare the steady-state error to a unit ramp input for (a) open-loop control with $K = 1$, and (b) closed-loop control with proportional controller $K = 5$ and unity feedback.

\textbf{Solution:}

\textbf{(a) Open-Loop:}

For open-loop control, the output is $Y(s) = G(s) K R(s)$ and the error is $E(s) = R(s) - Y(s)$.

For a unit ramp, $R(s) = 1/s^2$:
\begin{equation*}
Y(s) = \frac{10}{s(s+2)} \cdot 1 \cdot \frac{1}{s^2} = \frac{10}{s^3(s+2)}
\end{equation*}

The error is:
\begin{equation*}
E(s) = R(s) - Y(s) = \frac{1}{s^2} - \frac{10}{s^3(s+2)} = \frac{s(s+2) - 10}{s^3(s+2)}
\end{equation*}

Applying the Final Value Theorem:
\begin{equation*}
e_{ss} = \lim_{s \to 0} s \cdot E(s) = \lim_{s \to 0} \frac{s(s+2) - 10}{s^2(s+2)} = \lim_{s \to 0} \frac{s^2 + 2s - 10}{s^2(s+2)}
\end{equation*}

As $s \to 0$, the numerator approaches $-10$ while the denominator approaches $0$, indicating the error grows without bound. The open-loop system cannot track a ramp input.

\textbf{Alternatively}, since the open-loop system $GK$ is Type 1 (one integrator), but tracking requires matching the input, and the system does not inherently invert the input dynamics, ramp tracking fails.

$\boxed{e_{ss,\text{open-loop}} = \infty}$

\textbf{(b) Closed-Loop:}

The loop transfer function is:
\begin{equation*}
L(s) = G(s) K = \frac{10 \cdot 5}{s(s+2)} = \frac{50}{s(s+2)}
\end{equation*}

The velocity constant is:
\begin{equation*}
K_v = \lim_{s \to 0} s \cdot L(s) = \lim_{s \to 0} s \cdot \frac{50}{s(s+2)} = \lim_{s \to 0} \frac{50}{s+2} = \frac{50}{2} = 25
\end{equation*}

For a Type 1 system with ramp input:
\begin{equation*}
e_{ss} = \frac{1}{K_v} = \frac{1}{25} = 0.04
\end{equation*}

$\boxed{e_{ss,\text{closed-loop}} = 0.04 \text{ (or 4\% of ramp slope)}}$

\textbf{Conclusion:} Feedback reduces the error from unbounded to a finite, small value.
\end{example}

\begin{examplebox}[Sensitivity Improvement Through Feedback]
This example demonstrates the power of feedback for reducing sensitivity to plant variations. A temperature control system has nominal plant gain $G_0 = 5$ °C/V, but due to environmental variations, the actual gain can vary by $\pm 20\%$. We will design a proportional feedback controller to reduce the output sensitivity to plant variations by a factor of 10, showing how feedback can make a system robust to uncertainty.
\end{examplebox}

\begin{example}[Sensitivity Improvement]
\label{ex:sensitivity}
A temperature control system has nominal plant gain $G_0 = 5$ °C/V. Due to environmental variations, the actual gain can vary by $\pm 20\%$, i.e., $G = G_0(1 + \Delta)$ where $|\Delta| \leq 0.2$.

Design a proportional feedback controller to reduce the output sensitivity to plant variations by a factor of 10.

\textbf{Solution:}

For an open-loop system, a 20\% plant variation causes a 20\% output variation. We want to reduce this to 2\%.

The sensitivity function for closed-loop control is:
\begin{equation*}
S = \frac{1}{1 + L} = \frac{1}{1 + GKH}
\end{equation*}

For unity feedback ($H = 1$) and nominal plant:
\begin{equation*}
S_0 = \frac{1}{1 + G_0 K}
\end{equation*}

We require $|S_0| = 0.1$ (factor of 10 reduction):
\begin{equation*}
\frac{1}{1 + G_0 K} = 0.1
\end{equation*}
\begin{equation*}
1 + G_0 K = 10
\end{equation*}
\begin{equation*}
G_0 K = 9
\end{equation*}
\begin{equation*}
K = \frac{9}{G_0} = \frac{9}{5} = 1.8 \text{ V/°C}
\end{equation*}

\textbf{Verification:}

With $K = 1.8$ and $G_0 = 5$, the loop gain is $L_0 = 9$.

If the plant varies by $+20\%$: $G = 6$, $L = 6 \times 1.8 = 10.8$
\begin{equation*}
T = \frac{L}{1+L} = \frac{10.8}{11.8} = 0.915
\end{equation*}

Nominal closed-loop gain: $T_0 = \frac{9}{10} = 0.9$

Variation in $T$: $\dfrac{0.915 - 0.9}{0.9} = 1.67\%$

A 20\% plant variation now causes only 1.67\% output variation, confirming the sensitivity reduction.

$\boxed{K = 1.8 \text{ V/°C}}$
\end{example}

\begin{example}[Disturbance Rejection Analysis]
\label{ex:disturbance}
Consider the motor control system with block diagram shown in Figure~\ref{fig:closedloop_block}. The plant is $G(s) = \dfrac{1}{s+1}$, the controller is $K(s) = K_p$ (proportional), and feedback is unity ($H = 1$). A constant disturbance $D(s) = D_0/s$ acts at the plant input. Determine (a) the steady-state output error due to the disturbance for the open-loop system, (b) the steady-state output error due to the disturbance for the closed-loop system with $K_p = 10$, and (c) what value of $K_p$ reduces the disturbance effect by a factor of 100.

\textbf{Solution:}

\textbf{(a) Open-Loop Disturbance Response:}

For open-loop control, the disturbance-to-output transfer function is simply $G(s)$:
\begin{equation*}
Y_d(s) = G(s) D(s) = \frac{1}{s+1} \cdot \frac{D_0}{s}
\end{equation*}

Steady-state (Final Value Theorem):
\begin{equation*}
y_{d,ss} = \lim_{s \to 0} s \cdot Y_d(s) = \lim_{s \to 0} \frac{D_0}{s+1} = D_0
\end{equation*}

$\boxed{y_{d,ss}^{\text{OL}} = D_0}$ (the full disturbance appears at output)

\textbf{(b) Closed-Loop Disturbance Response:}

From equation (\ref{eq:disturbance_tf}):
\begin{equation*}
\frac{Y(s)}{D(s)} = \frac{G(s)}{1 + G(s)K_p} = \frac{\frac{1}{s+1}}{1 + \frac{K_p}{s+1}} = \frac{1}{s+1+K_p}
\end{equation*}

For step disturbance:
\begin{equation*}
Y_d(s) = \frac{1}{s+1+K_p} \cdot \frac{D_0}{s}
\end{equation*}

Steady-state:
\begin{equation*}
y_{d,ss} = \lim_{s \to 0} s \cdot Y_d(s) = \frac{D_0}{1+K_p} = \frac{D_0}{1+10} = \frac{D_0}{11}
\end{equation*}

$\boxed{y_{d,ss}^{\text{CL}} = D_0/11 \approx 0.091 D_0}$ (disturbance attenuated by factor of 11)

\textbf{(c) Design for 100$\times$ Attenuation:}

We require:
\begin{equation*}
\frac{y_{d,ss}^{\text{CL}}}{y_{d,ss}^{\text{OL}}} = \frac{1}{100}
\end{equation*}
\begin{equation*}
\frac{D_0/(1+K_p)}{D_0} = \frac{1}{1+K_p} = \frac{1}{100}
\end{equation*}
\begin{equation*}
1 + K_p = 100
\end{equation*}

$\boxed{K_p = 99}$
\end{example}

\begin{example}[Complete Closed-Loop Design]
\label{ex:design}
Design a feedback control system for a thermal system with the following specifications: the plant transfer function is $G(s) = \dfrac{2}{5s + 1}$ (°C per watt of heater power), the steady-state error to a step change in setpoint must be $\leq 2\%$, and the sensor has transfer function $H(s) = 1$ (ideal thermocouple). Determine the required proportional controller gain $K_p$.

\textbf{Solution:}

For a unity feedback system with proportional controller, the closed-loop transfer function is:
\begin{equation*}
T(s) = \frac{G(s) K_p}{1 + G(s) K_p} = \frac{\frac{2K_p}{5s+1}}{1 + \frac{2K_p}{5s+1}} = \frac{2K_p}{5s + 1 + 2K_p}
\end{equation*}

The DC gain (setting $s = 0$):
\begin{equation*}
T(0) = \frac{2K_p}{1 + 2K_p}
\end{equation*}

For a step input of magnitude $R_0$, the steady-state output is:
\begin{equation*}
y_{ss} = T(0) \cdot R_0 = \frac{2K_p}{1 + 2K_p} R_0
\end{equation*}

The steady-state error is:
\begin{equation*}
e_{ss} = R_0 - y_{ss} = R_0 \left(1 - \frac{2K_p}{1 + 2K_p}\right) = R_0 \cdot \frac{1}{1 + 2K_p}
\end{equation*}

The fractional error is:
\begin{equation*}
\frac{e_{ss}}{R_0} = \frac{1}{1 + 2K_p} \leq 0.02
\end{equation*}

Solving:
\begin{equation*}
1 + 2K_p \geq 50
\end{equation*}
\begin{equation*}
K_p \geq 24.5
\end{equation*}

$\boxed{K_p \geq 24.5 \text{ W/°C}}$

\textbf{Verification:} With $K_p = 25$:
\begin{equation*}
e_{ss}/R_0 = \frac{1}{1 + 50} = \frac{1}{51} \approx 1.96\% < 2\% \quad \checkmark
\end{equation*}

\textbf{Note:} This analysis assumes the closed-loop system is stable, which should be verified. The characteristic equation is $5s + 1 + 2K_p = 0$, giving pole at $s = -(1 + 2K_p)/5$. For any $K_p > 0$, this pole is in the left half-plane, confirming stability.
\end{example}

%=============================================================================
\section{Figures for TikZ Implementation}
\label{sec:tikz}
%=============================================================================

The following detailed descriptions enable creation of publication-quality figures using TikZ.

\subsection{Figure: Watt Centrifugal Governor}

\begin{figure}[htbp]
\centering
\begin{tikzpicture}[>=latex', scale=0.9]
    % Spindle
    \draw[thick] (0,0) -- (0,5);
    \draw[thick] (0,0) -- (0,-0.5);
    \fill (0,5) circle (2pt);

    % Low speed configuration (left side, dashed)
    \draw[dashed, gray] (0,5) -- (-1.2,3.5);
    \fill[gray] (-1.2,3.5) circle (4pt);
    \draw[dashed, gray] (-1.2,3.5) -- (-0.3,2.5);

    % High speed configuration (solid)
    \draw[thick] (0,5) -- (-2,2.5);
    \fill (-2,2.5) circle (6pt);
    \node[left] at (-2.2,2.5) {Flyball};
    \draw[thick] (-2,2.5) -- (-0.3,1.5);

    % Mirror for right side
    \draw[thick] (0,5) -- (2,2.5);
    \fill (2,2.5) circle (6pt);
    \draw[thick] (2,2.5) -- (0.3,1.5);

    % Collar
    \draw[thick, fill=gray!30] (-0.3,1.3) rectangle (0.3,1.7);
    \node[right] at (0.5,1.5) {Collar};

    % Throttle linkage
    \draw[thick] (0.3,1.5) -- (2.5,1.5);

    % Throttle valve
    \draw[thick] (2.5,0.5) -- (2.5,2.5);
    \draw[thick, fill=gray!50] (2.3,1.3) rectangle (2.7,1.7);
    \draw[thick] (2.5,2.7) -- (2.5,3.2);
    \draw[thick] (2.5,0.3) -- (2.5,-0.2);
    \node[right] at (2.8,1.5) {Valve};

    % Steam flow arrows
    \draw[->, thick] (2.5,-0.5) -- (2.5,0.3);
    \node[below] at (2.5,-0.5) {\small From boiler};
    \draw[->, thick] (2.5,3.2) -- (2.5,4);
    \node[above] at (2.5,4) {\small To engine};

    % Drive connection
    \draw[thick] (-0.3,-0.5) -- (0.3,-0.5) -- (0.3,-0.3) -- (-0.3,-0.3) -- cycle;
    \node[below] at (0,-0.7) {\small Drive from engine};

    % Rotation arrow
    \draw[->, thick] (0.5,4.5) arc (0:270:0.3);
    \node[right] at (0.6,4.3) {$\omega$};

    % Force annotations
    \draw[->, thick, blue] (-2,2.5) -- (-2.8,2.5);
    \node[above, blue] at (-2.6,2.6) {\small $F_c$};
    \draw[->, thick, red] (-2,2.5) -- (-2,1.8);
    \node[left, red] at (-2.1,2.1) {\small $mg$};

    % Legend
    \node[align=left] at (-1.5,-1.2) {\small Dashed: low speed\\Solid: high speed};
\end{tikzpicture}
\caption{Watt centrifugal governor schematic. As rotational speed $\omega$ increases, centrifugal force $F_c$ swings the flyballs outward, raising the collar and closing the throttle valve to reduce steam flow.}
\label{fig:watt_governor}
\end{figure}

\subsection{Figure: Laboratory Setup}

\begin{figure}[htbp]
\centering
\begin{tikzpicture}[>=latex', scale=0.85,
    component/.style={rectangle, draw, minimum width=2cm, minimum height=1.2cm, align=center}]

    % Arduino
    \node[component, fill=blue!20] (arduino) at (0,0) {Arduino\\Uno};
    \node[below=0.1cm of arduino, font=\scriptsize] {D2,D3,D5,D6,A0};

    % Motor Driver
    \node[component, fill=green!20] (driver) at (4,0) {Motor\\Driver};
    \node[below=0.1cm of driver, font=\scriptsize] {L298N};

    % Motor with Encoder
    \node[component, fill=gray!30, minimum width=2.5cm] (motor) at (8,0) {DC Motor\\+ Encoder};

    % Potentiometer
    \node[component, fill=yellow!30, minimum width=1.5cm, minimum height=0.8cm] (pot) at (0,-2.5) {Pot};

    % Power Supply
    \node[component, fill=red!20, minimum width=1.5cm, minimum height=0.8cm] (power) at (4,2) {12V\\Supply};

    % Position indicator on motor
    \draw[thick] (9.5,0) -- (10,0);
    \draw[thick, ->] (10,0) -- (10.5,0.3);
    \node[right, font=\scriptsize] at (10.5,0.3) {Position};

    % Wiring connections
    % Arduino to Driver (control signals - green)
    \draw[->, green!60!black, thick] (arduino.east) -- node[above, font=\scriptsize] {D5,D6} (driver.west);

    % Driver to Motor (power - thick black)
    \draw[->, thick] (driver.east) -- node[above, font=\scriptsize] {OUT} (motor.west);

    % Encoder to Arduino (feedback - blue)
    \draw[->, blue, thick] (motor.south) -- (8,-1.5) -- (1,-1.5) -- node[left, font=\scriptsize] {D2,D3} (arduino.south east);

    % Power supply connections (red)
    \draw[->, red, thick] (power.south) -- node[right, font=\scriptsize] {12V} (driver.north);

    % Potentiometer to Arduino (yellow)
    \draw[->, yellow!70!black, thick] (pot.north) -- node[left, font=\scriptsize] {A0} (arduino.south);

    % Ground connections (black dashed)
    \draw[dashed] (arduino.south west) -- (-1,-1) -- (driver.south west);

    % Legend
    \node[align=left, font=\scriptsize] at (10,-2) {
        \textcolor{green!60!black}{---} Control\\
        \textcolor{blue}{---} Encoder\\
        \textcolor{red}{---} Power\\
        \textcolor{yellow!70!black}{---} Analog
    };
\end{tikzpicture}
\caption{Laboratory setup showing the Arduino microcontroller, motor driver, DC motor with encoder, potentiometer for setpoint input, and power supply. Arrows indicate signal flow direction.}
\label{fig:lab_setup}
\end{figure}

%=============================================================================
\section{Chapter Summary}
\label{sec:summary}
%=============================================================================

This chapter established the fundamental concepts of control systems. From a historical perspective, feedback control emerged from practical necessity---the need to regulate systems against unpredictable disturbances. From Ktesibios's water clock to Watt's steam governor, each invention solved problems that open-loop control could not address.

The distinction between open-loop and closed-loop systems is fundamental. Open-loop systems apply predetermined inputs without measuring outputs; they fail when disturbances occur or models are inaccurate. Closed-loop systems measure output error and adjust input to compensate, providing inherent robustness.

The mathematical framework developed in this chapter shows that the closed-loop transfer function $T(s) = \dfrac{GK}{1+GKH}$ differs fundamentally from the open-loop transfer function $GK$ through the $(1+L)$ denominator, which reduces sensitivity to plant variations and disturbances. Feedback reduces sensitivity to plant variations by the factor $S = 1/(1+L)$, and system type determines steady-state error characteristics for polynomial inputs.

The laboratory exercises demonstrate these principles using accessible hardware, showing quantitatively how closed-loop control improves accuracy and disturbance rejection compared to open-loop approaches.

\begin{tcolorbox}[colback=yellow!10!white,colframe=orange!75!black,title=Key Takeaway,sharp corners]
Feedback control enables systems to regulate themselves against disturbances and uncertainties that cannot be eliminated at their source. The fundamental principle---measure the error, apply correction---is universal across all control applications, from ancient water clocks to modern spacecraft.
\end{tcolorbox}

\newpage
%=============================================================================
\section{Exercises}
\label{sec:exercises}
%=============================================================================

\begin{exercise}
\textbf{(Historical Analysis)} Explain why Ktesibios's float valve constitutes a feedback controller. Identify the reference input, controlled variable, manipulated variable, and feedback element.
\end{exercise}

\begin{exercise}
\textbf{(Open vs.\ Closed Loop)} A coffee maker heats water for a fixed time (open-loop). Propose a closed-loop alternative and identify what sensor and actuator would be required.
\end{exercise}

\begin{exercise}
\textbf{(Steady-State Error)} For a system with $G(s) = \dfrac{5}{s+2}$, $K = 4$, and unity feedback, calculate the steady-state error to (a) a unit step input, and (b) a unit ramp input.
\end{exercise}

\begin{exercise}
\textbf{(Sensitivity Analysis)} A control system has loop gain $L_0 = 20$ at DC. If the plant gain increases by 10\%, what is the percentage change in closed-loop gain?
\end{exercise}

\begin{exercise}
\textbf{(Controller Design)} Design a proportional controller for $G(s) = \dfrac{3}{s+5}$ with unity feedback such that the steady-state error to a unit step is less than 5\%.
\end{exercise}

\begin{exercise}
\textbf{(Disturbance Rejection)} For the system in Exercise 5, calculate the steady-state output due to a unit step disturbance at the plant input, both with and without feedback.
\end{exercise}

\begin{exercise}
\textbf{(Laboratory Extension)} Modify the closed-loop Arduino code to implement proportional-integral (PI) control. What improvement do you expect in steady-state error?
\end{exercise}

\begin{exercise}
\textbf{(System Type)} Determine the system type and calculate all applicable error constants for the loop transfer function
\begin{equation*}
L(s) = \frac{100(s+2)}{s^2(s+5)(s+10)}
\end{equation*}
\end{exercise}

\begin{exercise}
\textbf{(Biological Feedback)} Human body temperature is regulated at approximately 37°C through sweating, shivering, and vasodilation/vasoconstriction. Draw a block diagram of this thermoregulatory system, identifying the setpoint, sensor, controller, and actuators. What type of controller action (proportional, integral, derivative, or combination) does the body appear to use?
\end{exercise}

\begin{exercise}
\textbf{(Closed-Loop Gain Calculation)} For the closed-loop system shown in Figure~\ref{fig:closedloop_block} with $G(s) = \dfrac{10}{s(s+5)}$, $K(s) = K_p$, and $H(s) = 1$:
\begin{enumerate}[(a)]
\item Derive the closed-loop transfer function $T(s) = Y(s)/R(s)$.
\item Find the value of $K_p$ that results in a closed-loop DC gain of 0.9.
\item Calculate the closed-loop poles for this value of $K_p$ and determine if the system is stable.
\end{enumerate}
\end{exercise}

%=============================================================================
% End of Chapter
%=============================================================================
